\documentclass[italian]{../../unipd-notes}

\unipdsetup{
  course = {Progettazione di Sistemi Concorrenti},
  short-course = {Sistemi Concorrenti},
  professor = {Prof.ssa Ada Lovelace},
  academic-year = {2026--2027},
  degree-year = {1},
  semester = {1},
  document-type = {Appunti ed esercitazioni di laboratorio},
  author = {Simone Siega},
  date = {3 agosto 2026},
  version = {1.0.0}
}


\begin{document}
\makecoursefrontmatter
\makecoursemainmatter

\chapter{Modello di un buffer limitato}

Un buffer limitato separa la velocità dei produttori da quella dei consumatori senza consentire una crescita incontrollata della memoria. Il modello adottato contiene al massimo $N$ elementi, protegge le transizioni di stato con un mutex e usa due semafori contatori per rappresentare le posizioni occupate e quelle libere. La stessa struttura compare nelle code di rete, nei driver e nelle pipeline di elaborazione.

\begin{definition}
Lo stato del buffer è la tripla $(h,t,n)$: $h$ indica la prossima posizione da leggere, $t$ la prossima posizione da scrivere e $n$ il numero di elementi presenti. Gli indici sono calcolati modulo la capacità $N$.
\end{definition}

\begin{theorem}\label{thm:bounds}
Ogni esecuzione che inizializza $n=0$, riserva una posizione libera prima di ogni inserimento, riserva una posizione occupata prima di ogni rimozione e modifica il contatore soltanto mentre il mutex è acquisito conserva l'invariante
\[
  0 \leq n \leq N.
\]
\end{theorem}
\begin{proof}
Un inserimento può iniziare solo dopo aver riservato una posizione libera; vale quindi $n<N$ e l'aggiornamento produce $n+1\leq N$. Una rimozione richiede invece un elemento disponibile, parte da $n>0$ e produce $n-1\geq0$. Il mutex ordina totalmente gli aggiornamenti e permette di concludere per induzione sul numero di operazioni.
\end{proof}

Per questa prova di impaginazione si suppone che una sezione critica non contesa richieda \qty{1.5}{\micro\second} e si usano i campioni illustrativi \numlist{8.2;8.7;9.1}\unit{\micro\second}. Questi valori sintetici mostrano la formattazione di quantità ed elenchi e non sono misure sperimentali.

\begin{theorem}\label{thm:fifo}
Con la disciplina di permessi e mutex della \cref{thm:bounds}, se l'inserimento scrive nella posizione $t$ prima di incrementarla e la rimozione legge la posizione $h$ prima di avanzare, gli elementi sono restituiti in ordine FIFO.
\end{theorem}
\begin{proof}
Entrambi gli indici avanzano di una posizione modulo $N$. Una cella non può essere riutilizzata prima della sua lettura, perché la disciplina dei permessi impedisce al produttore di superare il consumatore. Le rimozioni visitano pertanto le celle nello stesso ordine in cui sono state riempite.
\end{proof}

Le proprietà di sicurezza non garantiscono da sole il progresso. Uno scheduler non equo può rimandare indefinitamente un thread pronto; la latenza reale dipende quindi anche dall'implementazione dei semafori e dalle politiche del sistema operativo.

\chapter{Architettura e misure}

La \cref{fig:architecture} riassume il percorso dei dati. I produttori completano il parsing prima di entrare nella sezione critica; i consumatori liberano la posizione occupata prima di elaborare il messaggio. Il buffer assume la proprietà dell'elemento al termine dell'inserimento e la trasferisce nuovamente alla rimozione.

\begin{figure}[H]
  \centering
  \begin{tikzpicture}[node distance=16mm]
    \node[draw=unipd-primary,rounded corners] (producer) {Produttori};
    \node[draw=unipd-accent,rounded corners,right=of producer] (buffer) {Buffer limitato};
    \node[draw=unipd-primary,rounded corners,right=of buffer] (consumer) {Consumatori};
    \draw[draw=unipd-accent,-{Latex},thick] (producer) -- node[above]{messaggi} (buffer);
    \draw[draw=unipd-accent,-{Latex},thick] (buffer) -- node[above]{messaggi} (consumer);
  \end{tikzpicture}
  \caption{Trasferimento della proprietà attraverso il buffer}\label{fig:architecture}
\end{figure}

\begin{figure}[H]
  \centering
  \begin{tikzpicture}[x=1.3cm,y=0.7cm]
    \draw[->] (0,0) -- (7,0) node[right]{tempo};
    \draw[fill=unipd-soft,draw=unipd-accent] (0.4,0.3) rectangle (2.2,0.9);
    \draw[fill=unipd-accent!20,draw=unipd-accent] (2.2,0.3) rectangle (3.0,0.9);
    \draw[fill=unipd-soft,draw=unipd-accent] (3.0,0.3) rectangle (5.7,0.9);
    \node[anchor=base] at (1.3,0.5) {preparazione};
    \node[anchor=base] at (2.6,0.5) {mutex};
    \node[anchor=base] at (4.35,0.5) {elaborazione};
  \end{tikzpicture}
  \caption{Collocazione della sezione critica in una transazione}\label{fig:timeline}
\end{figure}

Il benchmark illustrativo seguente ipotizza un milione di messaggi dopo una fase di riscaldamento. Ogni valore sintetico rappresenta una possibile mediana di sette esecuzioni e il throughput è definito sulle rimozioni completate, non sui tentativi di inserimento. La \cref{tab:throughput} serve a verificare l'impaginazione della tabella e non descrive un esperimento.

\begin{table}[H]
  \centering
  \caption{Valori sintetici di throughput per diverse capacità del buffer}\label{tab:throughput}
  \begin{tabular}{rrr}
    \toprule
    Capacità & Throughput (messaggi/s) & Latenza al 99° percentile \\
    \midrule
    1 & 79000 & \qty{44}{\micro\second} \\
    16 & 236000 & \qty{20}{\micro\second} \\
    64 & 251000 & \qty{21}{\micro\second} \\
    256 & 249000 & \qty{29}{\micro\second} \\
    \bottomrule
  \end{tabular}
  \tablesource{dati sintetici creati per questo esempio di integrazione}
\end{table}

\begin{table}[H]
  \centering
  \caption{Transizioni di stato delle operazioni}\label{tab:transitions}
  \begin{tabularx}{\textwidth}{lLL}
    \toprule
    Operazione & Precondizione & Modifica atomica dello stato \\
    \midrule
    Inserimento & $n<N$ & $a[t]\gets x$, $t\gets(t+1)\bmod N$, $n\gets n+1$ \\
    Rimozione & $n>0$ & $x\gets a[h]$, $h\gets(h+1)\bmod N$, $n\gets n-1$ \\
    Ispezione & Mutex acquisito & Lettura di $(h,t,n)$ senza modifiche \\
    \bottomrule
  \end{tabularx}
\end{table}

\begin{unipdcode}[language=Python,caption={Inserimento con controllo esplicito},label={lst:insert}]
def enqueue(queue, item):
    with queue.mutex:
        if queue.count == queue.capacity:
            raise BufferError("buffer pieno")
        queue.data[queue.tail] = item
        queue.tail = (queue.tail + 1) % queue.capacity
        queue.count += 1
\end{unipdcode}

\chapter{Scheduling e verifica}

Un inserimento bloccante riserva prima una posizione libera e acquisisce poi il mutex. Invertire le due operazioni può causare un deadlock: il produttore conserverebbe il mutex in attesa di un consumatore che necessita dello stesso mutex. La rimozione usata nei test è riportata nel \cref{lst:remove}.

\begin{unipdcode}[language=Python,caption={Rimozione con controllo esplicito},label={lst:remove}]
def dequeue(queue):
    with queue.mutex:
        if queue.count == 0:
            raise BufferError("buffer vuoto")
        item = queue.data[queue.head]
        queue.head = (queue.head + 1) % queue.capacity
        queue.count -= 1
        return item
\end{unipdcode}

I controlli nei \cref{lst:insert,lst:remove} trasformano una violazione dell'invariante in un errore immediato. Il codice di produzione attende normalmente un semaforo, ma le asserzioni sul contatore restano utili per individuare segnali errati e accessi eseguiti senza lock. L'\cref{alg:insert} adotta l'ordine di acquisizione sicuro.

\begin{algorithm}[H]
  \caption{Inserimento bloccante}\label{alg:insert}
  \KwInput{Una coda $q$ e un elemento $x$}
  \KwOutput{La coda aggiornata}
  attendi($q.libere$)\;
  acquisisci($q.mutex$)\;
  $q.dati[q.coda]\gets x$\;
  $q.coda\gets(q.coda+1)\bmod q.capacita$\;
  rilascia($q.mutex$)\;
  segnala($q.usate$)\;
  \Return{$q$}\;
\end{algorithm}

La procedura nell'\cref{alg:validate} non modifica lo stato. Può essere eseguita dopo ogni operazione casuale di un test basato su modello, confrontando la coda reale con una semplice sequenza di riferimento.

\begin{algorithm}[H]
  \caption{Validazione dello stato della coda}\label{alg:validate}
  \KwInput{Lo stato di una coda $q$}
  \KwOutput{Vero se e solo se lo stato è valido}
  \If{$q.conteggio<0$ \KwOr $q.conteggio>q.capacita$}{
    \Return{falso}\;
  }
  \For{$i\gets 1$ \KwTo $q.conteggio$}{
    controlla l'elemento in $(q.testa+i-1)\bmod q.capacita$\;
  }
  \Return{vero}\;
\end{algorithm}

I riferimenti singolari minuscoli sono \cref{thm:bounds}, \cref{fig:architecture}, \cref{tab:throughput}, \cref{lst:insert} e \cref{alg:insert}. Le forme singolari maiuscole sono \Cref{thm:bounds}, \Cref{fig:architecture}, \Cref{tab:throughput}, \Cref{lst:insert} e \Cref{alg:insert}.

Le forme plurali minuscole sono \cref{thm:bounds,thm:fifo}, \cref{fig:architecture,fig:timeline}, \cref{tab:throughput,tab:transitions}, \cref{lst:insert,lst:remove} e \cref{alg:insert,alg:validate}. Le forme plurali maiuscole sono \Cref{thm:bounds,thm:fifo}, \Cref{fig:architecture,fig:timeline}, \Cref{tab:throughput,tab:transitions}, \Cref{lst:insert,lst:remove} e \Cref{alg:insert,alg:validate}.

\printcourselists

\courseappendices
\chapter{Lista di controllo per il laboratorio}

Prima delle misure, compilare il programma con le ottimizzazioni attive ma senza eliminare le asserzioni. È opportuno fissare l'affinità dei thread quando si confrontano strategie diverse, annotare il modello del processore e la versione del sistema operativo, e mantenere invariato il carico di ingresso. Una fase iniziale non misurata stabilizza cache e frequenza della CPU.

Dopo ogni operazione si deve verificare $0\leq n\leq N$ e confrontare la sequenza rimossa con il modello di riferimento. Le capacità $1$, $2$ e un valore non potenza di due sono casi particolarmente efficaci: il primo espone errori di ordinamento, mentre gli altri sollecitano il riavvolgimento degli indici.

La relazione finale riporta revisione del sorgente, opzioni del compilatore, numero di ripetizioni, metodo di aggregazione e osservazioni grezze. Throughput e latenza devono rimanere distinti; gli outlier vanno spiegati, non eliminati senza motivazione. In questo modo l'esperimento resta riproducibile e le conclusioni sono legate all'implementazione effettivamente misurata.

\end{document}
